% =============================================================
% Bilgisayar Mimarisi Kitabı - Preamble
% Oğuz Ergin
% Şablon C: Modern Teknik (Libertinus + renkli başlıklar)
% =============================================================

% --- Temel Paketler ---
\usepackage{fontspec}
\usepackage[turkish,shorthands=:!]{babel}

% --- Yazı Tipi: Libertinus ---
\usepackage{libertinus}
\setmonofont{Consolas}[Scale=0.82]    % Monospace: Consolas (Unicode/Türkçe destekli)

% --- Sayfa Düzeni ---
\usepackage[a4paper, margin=2.5cm, top=3cm, bottom=3cm, headheight=25pt]{geometry}

% --- Renkler (font/listing/titlesec tanımlarından önce) ---
\usepackage{xcolor}
\usepackage{pifont}  % \ding{51} (checkmark) ve \ding{55} (cross) için
\definecolor{kuturenk}{RGB}{230, 240, 250}
\definecolor{baslikrenk}{RGB}{0, 71, 133}
\definecolor{bolumrenk}{RGB}{0, 71, 133}
\definecolor{acikgri}{RGB}{245, 245, 245}

% TikZ blok şeması renkleri tikz-sekiller-makrolar.tex'te tanımlıdır.

% --- Bölüm Başlık Stili ---
\usepackage{titlesec}

% Bölüm başlığı: Renkli büyük numara + alt çizgi
\titleformat{\chapter}[display]
    {\normalfont\bfseries\color{bolumrenk}}
    {\flushright\fontsize{72}{72}\selectfont\color{bolumrenk!30}\thechapter}
    {-20pt}
    {\huge\flushleft}
    [\vspace{-5pt}{\color{bolumrenk}\titlerule[2pt]}]

% Section: renkli
\titleformat{\section}
    {\normalfont\Large\bfseries\color{bolumrenk}}
    {\thesection}{1em}{}

% Subsection
\titleformat{\subsection}
    {\normalfont\large\bfseries\color{bolumrenk!80}}
    {\thesubsection}{1em}{}

% --- Matematik ---
\usepackage{amsmath}
\usepackage{unicode-math}
\setmathfont{Libertinus Math}

% Checkmark: Libertinus'ta ✓ glifi yok, platformda bulunan fonttan al
\IfFontExistsTF{Apple Symbols}{%
  \newfontfamily\symbolfont{Apple Symbols}%
}{%
  \newfontfamily\symbolfont{Segoe UI Symbol}%
}
\AtBeginDocument{\def\checkmark{{\symbolfont ✓}}}

% --- Tablolar ---
\usepackage{booktabs}
\usepackage{array}
\usepackage{multirow}
\usepackage{longtable}
\usepackage{tabularx}

% --- Listeler ---
\usepackage[shortlabels]{enumitem}

% --- Şekiller ve Grafikler ---
\usepackage{graphicx}
\usepackage{float}
\usepackage{subcaption}
\captionsetup{labelfont=bf}                    % "Tablo 1.1:" ve "Şekil 1.1:" kalın
\captionsetup[table]{textfont=it}              % Tablo açıklama metni yatık
\usepackage{tikz}
\usetikzlibrary{shapes, arrows, arrows.meta, positioning, calc, decorations.markings,
    circuits.logic.US, circuits.logic.IEC, fit, patterns, matrix, automata, babel,
    decorations.pathreplacing, decorations.pathmorphing, shadows, backgrounds}
\usepackage[section]{placeins}  % Şekillerin bölüm dışına taşmasını engeller

% TikZ global ok stili — tüm tikzpicture'larda >=Stealth varsayılan olur
\tikzset{>=Stealth}
\tikzset{okstil/.style={->, >=Stealth, thick}}

% --- TikZ Şekil Makroları ---
\input{tikz-sekiller-makrolar}

% --- Kod Listeleme ---
\usepackage{listings}
\lstloadlanguages{[ANSI]C}  % listings 1.11b uyumsuzluk düzeltmesi: C dilini önceden yükle
\lstset{
    basicstyle=\ttfamily\footnotesize,
    breaklines=true,
    frame=single,
    numbers=left,
    numberstyle=\tiny\color{gray},
    tabsize=4,
    captionpos=b,
    xleftmargin=2em,
    framexleftmargin=1.5em,
    backgroundcolor=\color{acikgri},
    rulecolor=\color{bolumrenk!40},
    extendedchars=true,
    columns=fullflexible,
    keepspaces=true,
}

% --- Assembly/MIPS stil tanımı ---
\lstdefinestyle{mips}{
    morekeywords={add, sub, addi, lw, sw, beq, bne, slt, j, jal, jr, and, or, sll, srl, lui, ori, mult, div, mfhi, mflo},
    keywordstyle=\color{blue}\bfseries,
    commentstyle=\color{green!50!black}\ttfamily,
    stringstyle=\color{red},
    morecomment=[l]{\#},
    texcl=true,
}

% --- Assembly/RISC-V stil tanımı ---
\lstdefinestyle{riscv}{
    morekeywords={add, sub, addi, and, or, xor, andi, ori, xori, sll, srl, sra, slli, srli, srai, slt, sltu, slti, sltiu, lw, sw, lb, sb, lh, sh, lbu, lhu, ld, sd, beq, bne, blt, bge, bltu, bgeu, jal, jalr, lui, auipc, ecall, ebreak, mul, mulh, mulhu, mulhsu, div, divu, rem, remu, nop, li, la, mv, not, neg, j, ret, call, tail, fadd.s, fsub.s, fmul.s, fdiv.s, fsqrt.s, fadd.d, fsub.d, fmul.d, fdiv.d, flw, fsw, fld, fsd, feq.s, flt.s, fle.s, fcvt.w.s, fcvt.s.w, csrr, csrw, csrs, csrc, fence, fence.i, wfi, mret, sret, lr.w, sc.w, lr.d, sc.d, amoswap.w, amoadd.w, amoand.w, amoor.w, amoxor.w, amomax.w, amomin.w, amomaxu.w, amominu.w, amoswap.d, amoadd.d, amoand.d, amoor.d, amoxor.d, amomax.d, amomin.d, amomaxu.d, amominu.d, prefetch.r, prefetch.w, prefetch.i, zero, ra, sp, gp, tp, t0, t1, t2, t3, t4, t5, t6, s0, s1, s2, s3, s4, s5, s6, s7, s8, s9, s10, s11, a0, a1, a2, a3, a4, a5, a6, a7, fp, x0, x1, x2, x3, x4, x5, x6, x7, x8, x9, x10, x11, x12, x13, x14, x15, x16, x17, x18, x19, x20, x21, x22, x23, x24, x25, x26, x27, x28, x29, x30, x31, f0, f1, f2, f3, f4, f5, f6, f7},
    keywordstyle=\color{blue}\bfseries,
    commentstyle=\color{green!50!black}\ttfamily,
    stringstyle=\color{red},
    morecomment=[l]{\#},
    texcl=true,
    sensitive=true,
}

% --- Verilog stil tanımı ---
\lstdefinestyle{verilog}{
    morekeywords={module, endmodule, input, output, inout, wire, reg, parameter, localparam, assign, always, begin, end, if, else, case, endcase, default, posedge, negedge, initial, forever, for, while, integer, genvar, generate, endgenerate, task, endtask, function, endfunction},
    keywordstyle=\color{blue}\bfseries,
    commentstyle=\color{green!50!black},
    stringstyle=\color{red},
    sensitive=true,
}

% --- CUDA/C stil tanımı ---
\lstdefinestyle{cuda}{
    morekeywords={__global__, __device__, __shared__, __host__, __constant__, int, void, if, else, return, for, while, float, double, const, unsigned, struct, typedef, sizeof, static, extern, char, short, long, signed, blockIdx, blockDim, threadIdx, gridDim, warpSize},
    keywordstyle=\color{blue}\bfseries,
    commentstyle=\color{green!50!black},
    stringstyle=\color{red},
}

% --- Alıştırmalardaki kod blokları için kompakt stil ---
\lstdefinestyle{alistirma-kod}{
    numbers=none,
    xleftmargin=1em,
    framexleftmargin=0.5em,
    frame=single,
    backgroundcolor=\color{acikgri},
    rulecolor=\color{bolumrenk!40},
    basicstyle=\ttfamily\small,
    breaklines=true,
    tabsize=4,
    aboveskip=0.5\baselineskip,
    belowskip=0.5\baselineskip,
}

% --- Bağlantılar ---
\usepackage{hyperref}
\hypersetup{
    colorlinks=true,
    linkcolor=bolumrenk,
    citecolor=bolumrenk,
    urlcolor=bolumrenk!70,
    bookmarks=true,
    bookmarksnumbered=true,
}

% --- Kaynakça ---
\usepackage{csquotes}
\DeclareQuoteAlias{american}{turkish}
\usepackage[style=numeric, sorting=nyt, backend=biber, defernumbers=true]{biblatex}
\addbibresource{kaynakca.bib}

% --- Dizin ---
\usepackage{makeidx}
\makeindex

% --- Özel Ortamlar (Environments) ---
\usepackage{tcolorbox}
\tcbuselibrary{breakable, skins}

% Örnek kutusu
\newtcolorbox{ornek}[1][]{
    colback=bolumrenk!5,
    colframe=bolumrenk,
    fonttitle=\bfseries,
    title={Örnek #1},
    breakable,
    enhanced,
    left=4mm,
    right=4mm,
    arc=1mm,
}

% Tanım kutusu
\newtcolorbox{tanim}[1][]{
    colback=orange!5!white,
    colframe=orange!80!black,
    fonttitle=\bfseries,
    title={#1},
    breakable,
    enhanced,
    arc=1mm,
}

% Alıştırma ortamı
\newcounter{alistirma}[chapter]
\renewcommand{\thealistirma}{\thechapter.\arabic{alistirma}}
\newenvironment{alistirma}{%
    \refstepcounter{alistirma}%
    \par\medskip\noindent\textbf{Alıştırma \thealistirma.}\space
}{%
    \par\medskip
}

% Bilgi notu sayacı ve listesi
\newcounter{bilginotu}[chapter]
\renewcommand{\thebilginotu}{\thechapter.\arabic{bilginotu}}

\makeatletter
\newcommand{\listofbilginotlari}{%
    \chapter*{Bilgi Notları Listesi}%
    \addcontentsline{toc}{chapter}{Bilgi Notları Listesi}%
    \@starttoc{lon}%
}
\newcommand{\l@bilginotu}{\@dottedtocline{1}{0em}{3em}}
\makeatother

% Bilgi notu iç kutusu (doğrudan kullanılmaz)
\newtcolorbox{bilginotuic}[1]{
    colback=green!5!white,
    colframe=green!60!black,
    fonttitle=\bfseries,
    title={#1},
    enhanced,
    arc=1mm,
    borderline west={3pt}{0pt}{green!60!black},
}

% Bilgi notu ortamı (numaralı, listeye eklenir)
\newenvironment{bilginot}[1][]{%
    \refstepcounter{bilginotu}%
    \def\bilginotarg{#1}%
    \ifx\bilginotarg\empty
        \addcontentsline{lon}{bilginotu}{\protect\numberline{\thebilginotu}Bilgi Notu}%
        \begin{bilginotuic}{Bilgi Notu~\thebilginotu}%
    \else
        \addcontentsline{lon}{bilginotu}{\protect\numberline{\thebilginotu}#1}%
        \begin{bilginotuic}{Bilgi Notu~\thebilginotu: #1}%
    \fi
}{%
    \end{bilginotuic}%
}

% Terim kökeni kutusu
\definecolor{terimrenk}{RGB}{139, 90, 43}
\newtcolorbox{terimnotu}[1][Terim Kökenleri]{
    colback=terimrenk!5!white,
    colframe=terimrenk!80!black,
    fonttitle=\bfseries,
    title={\small\textsf{#1}},
    breakable,
    enhanced,
    arc=1mm,
    borderline west={3pt}{0pt}{terimrenk!80!black},
    fontupper=\small,
    top=2mm, bottom=2mm,
}

% --- Üstbilgi/Altbilgi ---
\usepackage{fancyhdr}
\pagestyle{fancy}
\fancyhf{}
\fancyhead[LE]{\small\color{bolumrenk}\leftmark}
\fancyhead[RO]{\small\color{bolumrenk}\rightmark}
\fancyfoot[C]{\color{bolumrenk}\thepage}
\renewcommand{\headrulewidth}{0.4pt}
\renewcommand{\headrule}{\hbox to\headwidth{\color{bolumrenk}\leaders\hrule height \headrulewidth\hfill}}
\renewcommand{\footrulewidth}{0pt}

% --- Şekil numaralama ---
\numberwithin{figure}{chapter}
\numberwithin{table}{chapter}
\numberwithin{equation}{chapter}

% --- Şekil/Tablo listesinde bölüm başlıkları ---
\makeatletter
\let\kitap@origchapter\@chapter
\def\@chapter[#1]#2{%
  \kitap@origchapter[#1]{#2}%
  \addtocontents{lof}{%
    {\protect\noindent\protect\bfseries\protect\color{bolumrenk}%
     \@chapapp~\thechapter{} --- #1\protect\par\protect\nopagebreak\protect\smallskip}%
  }%
  \addtocontents{lot}{%
    {\protect\noindent\protect\bfseries\protect\color{bolumrenk}%
     \@chapapp~\thechapter{} --- #1\protect\par\protect\nopagebreak\protect\smallskip}%
  }%
  \addtocontents{lon}{%
    {\protect\noindent\protect\bfseries\protect\color{bolumrenk}%
     \@chapapp~\thechapter{} --- #1\protect\par\protect\nopagebreak\protect\smallskip}%
  }%
}
\makeatother

% --- Bölüm açılış fotoğrafı ---
\newcommand{\bolumfoto}[2]{%
    \thispagestyle{empty}%
    \vspace*{\fill}%
    \begin{center}%
    \begin{minipage}{\textwidth}%
    \centering
    \includegraphics[width=\textwidth,height=0.6\textheight,keepaspectratio]{#1}\\[8pt]%
    {\small\color{bolumrenk}\textit{#2}}%
    \end{minipage}%
    \end{center}%
    \vspace*{\fill}%
    \newpage%
}

% --- Eksik şekil yer tutucusu ---
\newcommand{\eksikfigur}[2][]{%
    \begin{figure}[H]
    \centering
    \fbox{\parbox{0.7\textwidth}{\centering\vspace{2cm}\textbf{[ŞEKİL GEREKLİ]}\\#2\vspace{2cm}}}
    \caption{#2}
    \ifx&#1&\else\label{#1}\fi
    \end{figure}
}
